% Compiler avec XeLaTeX ou LuaLaTeX
\documentclass[11pt,a4paper]{article}

\usepackage{amsmath,amssymb}
\usepackage{fontspec}
\usepackage{unicode-math}
\usepackage{polyglossia}
\setmainlanguage{french}
\usepackage{xcolor}
\usepackage{fancyhdr}
\usepackage[margin=2cm,headheight=16pt,headsep=22pt]{geometry}
\usepackage{enumitem}

\IfFontExistsTF{Libertinus Serif}
  {\setmainfont{Libertinus Serif}}
  {\IfFontExistsTF{Latin Modern Roman}{\setmainfont{Latin Modern Roman}}{\setmainfont{lmroman10-regular.otf}}}

\IfFontExistsTF{Libertinus Math}
  {\setmathfont{Libertinus Math}}
  {\IfFontExistsTF{Latin Modern Math}{\setmathfont{Latin Modern Math}}{\setmathfont{latinmodern-math.otf}}}

\linespread{1.07}
\setlength{\parindent}{0pt}
\AtBeginDocument{\setlength{\parindent}{0pt}}
\setlength{\parskip}{0.35em}

\pagestyle{fancy}
\fancyhf{}
\lhead{Djamel BENMAKHLOUF}
\rhead{Terminale Spécialité Maths}
\usepackage{lastpage}
\cfoot{Page \thepage\ / \pageref{LastPage}}
\renewcommand{\headrulewidth}{0.5pt}
\renewcommand{\headrule}{\vskip 5pt\hrule width\headwidth height\headrulewidth \vskip-\headrulewidth}

\setlist[enumerate,1]{itemsep=0.45em, topsep=0.3em, label=\textbf{\arabic*.}}
\setlist[enumerate,2]{itemsep=0.25em, topsep=0.2em, label=\textbf{\alph*.}}

\newcommand{\exo}[2]{%
  \vspace{1.5em}%
  \noindent{\large\textbf{Exercice #1. -- #2}}%
  \par\vspace{0.35em}%
  \hrule height 0.5pt%
  \vspace{0.7em}%
}

\newcommand{\partie}[1]{%
  \vspace{0.9em}%
  \noindent\textbf{#1}%
  \par\vspace{0.3em}%
}

\begin{document}

\begin{center}
  {\LARGE \textbf{Présentation de NeoHarness AI}}\\[0.5em]
  {\large Terminale Spécialité Mathématiques}
\end{center}

\exo{1}{Conception et Architecture de NeoHarness AI}

\partie{Partie A -- Origine et Conception}

\begin{enumerate}
  \item Rappeler qui est le concepteur de NeoHarness AI.
  \item Expliquer le rôle du système d'orchestration et la distinction entre le concepteur et le modèle moteur utilisé.
\end{enumerate}

\partie{Partie B -- Correction détaillée}

\begin{enumerate}
  \item NeoHarness AI a été conçu et développé par Djamel BENMAKHLOUF, professeur de mathématiques, assisté par IA.
  \item Le harness orchestre différents modèles d'intelligence artificielle utilisés comme moteurs pour faire fonctionner l'agent, sans que le modèle moteur ne soit le concepteur du système.
\end{enumerate}

\end{document}
